% !TeX root = ../navier-stokes-manuscript.tex
\subsection{The Heat Clock and the Zeno Cascade}\label{sub:BodyHeatClock}

Under the Navier--Stokes scaling, a spatial length divided by \(\lambda\) comes with a time divided by \(\lambda^2\).
Viscosity carries this quadratic clock at every active scale.
To see it directly, let \(v_t=\nu\Delta v\) and suppose the Fourier mass of \(v\) lies near frequencies \(|\xi|\simeq r^{-1}\).
The heat flow gives
\[
  \widehat v(\xi,t+\delta t)
  =e^{-\nu|\xi|^2\delta t}\widehat v(\xi,t).
\]
An order-one viscous response occurs when \(\nu|\xi|^2\delta t\) is of order one.
At spatial scale \(r\), the corresponding heat time is
\begin{equation}\label{eq:BodyHeatComparisonTime}
  \tau_\nu(r)
  \asymp\frac{r^2}{\nu}.
\end{equation}
Halving the spatial scale quarters the heat time.

If \(T_*\) is a proposed terminal time, the spatial scale whose heat clock matches the remaining time satisfies
\[
  r_\nu(t)
  \asymp\sqrt{\nu(T_*-t)}.
\]
This is a comparison scale supplied by viscosity, not a claim that the solution has formed a vortex of that radius.

Now take a geometric sequence of shrinking scales,
\[
  r_k=2^{-k}r_0.
\]
Their heat clocks satisfy
\[
  \tau_k
  \asymp\frac{r_k^2}{\nu}
  =\frac{r_0^2}{\nu}4^{-k},
\]
and hence
\[
  \sum_{k=0}^{\infty}\tau_k
  \asymp
  \frac{r_0^2}{\nu}
  \sum_{k=0}^{\infty}4^{-k}
  <\infty.
\]
The heat-clock calculation permits successively smaller comparison scales to have a finite total time.
It does not show that one Navier--Stokes solution realizes those contractions.
The nonlinear evolution would still have to create each next scale from the same velocity field while diffusion is already acting on it.

The heat clock alone does not decide which quantity grows.
In the general critical-height route, the dangerous conditional is that the active scale falls while \(H(t)\) becomes unbounded.
Suppose \(H\) is unbounded along a finite maximal interval \([0,T_*)\).
The unbounded rise can be recorded by first passages.
Choose \(H_0>H(0)\), set
\[
  L_k:=2^kH_0,
  \qquad
  t_k:=\inf\{t>0:H(t)=L_k\}.
\]
Suppose the \(k\)-th first passage occurs while the relevant active scale is \(r_k\).
If the next crossing occurs within one fixed multiple of the heat time at that scale,
\[
  0<t_{k+1}-t_k
  \leq C\frac{r_k^2}{\nu},
\]
then the crossing times have a finite total:
\[
  \sum_{k=0}^{\infty}(t_{k+1}-t_k)<\infty,
  \qquad
  t_k\uparrow T_*<\infty.
\]
Under that hypothesis,
\[
  r_k\downarrow0,
  \qquad
  t_k\uparrow T_*<\infty,
  \qquad
  H(t_k)=L_k\uparrow\infty.
\]
The same velocity history passes through unbounded critical heights on shorter and shorter clocks.
This is the finite-time Zeno cascade.
These conclusions concern the critical height; they do not by themselves imply velocity-amplitude blow-up.
Spatial gradient growth likewise requires the velocity difference across the active width to decay more slowly than the width itself.

The height crossings do force a change in the \(\dot H^{1/2}\) velocity:
\[
  \norm{u(t_{k+1})-u(t_k)}_{\dot H^{1/2}}
  \geq
  \left|\sqrt{2L_{k+1}}-\sqrt{2L_k}\right|
  =(\sqrt2-1)\sqrt{2L_k}.
\]
Consequently,
\[
  \operatorname*{ess\,sup}_{t_k<t<t_{k+1}}
  \norm{\partial_tu(t)}_{\dot H^{1/2}}
  \geq
  \frac{(\sqrt2-1)\sqrt{2L_k}}{t_{k+1}-t_k}.
\]
Thus the Zeno crossings supply both shrinking intervals and growing Eulerian velocity responses.
The next question is whether later responses are activated in turn.
